% Part: computability
% Chapter: computability-theory
% Section: rice-theorem

\documentclass[../../../include/open-logic-section]{subfiles}

\begin{document}

\olfileid{cmp}{thy}{rce}
\olsection{قضیهٔ رایس}

اگر بیندیشید، می‌بینید ویژگی‌های خاص $\fn{Tot}$ نقشی در برهان
\olref[tot]{prop:total} ندارند. $h(x,y)$ را چنان ساختیم که دقیقاً
وقتی $x\in K$ مانند تابع ثابت $j(y)=0$ رفتار کند؛ اما در همان شرایط
می‌توانستیم به همان خوبی آن را شبیه هر تابع محاسبه‌پذیر جزئی دیگری
بسازیم. این مشاهده اجازه می‌دهد قضیه‌ای کلی‌تر بیان کنیم که به‌طور
تقریبی می‌گوید هیچ ویژگی نابدیهیِ توابع محاسبه‌پذیر تصمیم‌پذیر نیست.

به یاد داشته باشید که $\cfind{0}$، $\cfind{1}$، $\cfind{2}$،~\dots
برشماری معیار ما از توابع محاسبه‌پذیر جزئی است.

\begin{thm}[قضیهٔ رایس]
  بگذارید $C$ هر مجموعه‌ای از توابع محاسبه‌پذیر جزئی و
  $A=\Setabs{n}{\cfind{n}\in C}$ باشد. اگر $A$ محاسبه‌پذیر باشد،
  آنگاه یا $C$ تهی است یا $C$ مجموعهٔ همهٔ توابع محاسبه‌پذیر جزئی
  است.
\end{thm}

\emph{مجموعهٔ نمایه‌ها} مجموعه‌ای چون $A$ با این ویژگی است که اگر
$n$ و~$m$ نمایه‌هایی باشند که تابع یکسانی را «محاسبه» می‌کنند، آنگاه
یا هر دو در $A$ هستند یا هیچ‌یک نیست. به‌آسانی می‌توان دید مجموعهٔ
$A$ در قضیه این ویژگی را دارد. برعکس، اگر $A$ مجموعهٔ نمایه‌ها و $C$
مجموعهٔ توابع محاسبه‌شده با این نمایه‌ها باشد، آنگاه
$A=\Setabs{n}{\cfind{n}\in C}$.

\begin{explain}
با این اصطلاح، قضیهٔ رایس هم‌ارز این است که هیچ مجموعهٔ نمایه‌های
نابدیهی تصمیم‌پذیر نیست. برای فهم قضیه مفید است بر تمایز میان
\emph{برنامه‌ها} (مثلاً در زبان برنامه‌نویسی دلخواهتان) و توابعی که
محاسبه می‌کنند تأکید کنیم. مسلماً پرسش‌هایی محاسبه‌پذیر دربارهٔ
برنامه‌ها (نمایه‌ها)، که اشیایی نحوی‌اند، وجود دارد: آیا برنامه بیش
از ۱۵۰ نماد دارد؟ بیش از ۲۲ سطر دارد؟ دستور «while» دارد؟ آیا رشتهٔ
«hello world» جایی به‌عنوان آرگومان دستور «print» ظاهر می‌شود؟ قضیهٔ
رایس می‌گوید هیچ پرسش نابدیهی دربارهٔ \emph{رفتار} برنامه محاسبه‌پذیر
نیست. پرسش‌های زیر از این نوع‌اند: آیا برنامه روی ورودی~$0$ متوقف
می‌شود؟ آیا اصلاً متوقف می‌شود؟ آیا هیچ‌گاه عددی زوج بیرون می‌دهد؟
\end{explain}

\begin{proof}[برهان قضیهٔ رایس]
فرض کنید $C$ نه تهی باشد و نه مجموعهٔ همهٔ توابع محاسبه‌پذیر جزئی، و
$A$ مجموعهٔ نمایه‌های توابع در~$C$ باشد. نشان می‌دهیم اگر $A$
محاسبه‌پذیر بود، می‌توانستیم مسئلهٔ توقف را حل کنیم؛ پس $A$
محاسبه‌پذیر نیست.

بدون کاستن از کلیت می‌توان فرض کرد تابع $f$ که هیچ‌جا تعریف نشده است
در $C$ نیست (وگرنه در استدلال زیر $C$ و متممش را جابه‌جا کنید).
بگذارید $g$ هر تابعی در~$C$ باشد. اندیشه این است که اگر می‌توانستیم
دربارهٔ~$A$ تصمیم بگیریم، می‌توانستیم نمایه‌های محاسبه‌کنندهٔ~$f$ را
از نمایه‌های محاسبه‌کنندهٔ~$g$ بازشناسیم؛ و سپس از این توانایی برای
حل مسئلهٔ توقف استفاده کنیم.

چنین عمل می‌کنیم. با استفاده از توابع محاسبه‌پذیر جزئیِ جهانی، تابع
زیر را تعریف می‌کنیم:
\[
h(x,y) \simeq
\begin{cases}
\text{تعریف‌نشده} & \text{اگر $\cfind{x}(x) \fundefined$} \\
g(y) & \text{در غیر این صورت.}
\end{cases}
\]
برای محاسبهٔ $h$ نخست می‌کوشیم $\cfind{x}(x)$ را محاسبه کنیم؛ اگر آن
محاسبه متوقف شود، به محاسبهٔ~$g(y)$ می‌پردازیم؛ و اگر \emph{آن}
محاسبه متوقف شود، خروجی را بازمی‌گردانیم. به‌طور رسمی‌تر می‌توان نوشت:
\[
h(x,y) \simeq \Proj{2}{0}(g(y),\fn{Un}(x,x)).
\]
که در آن $\Proj{2}{0}(z_0,z_1)=z_0$ تابع افکنش دوموضعیِ
محاسبه‌پذیری است که آرگومان صفرم را بازمی‌گرداند.

پس $h$ ترکیبی از توابع محاسبه‌پذیر جزئی است و سمت راست دقیقاً وقتی
تعریف شده و برابر~$g(y)$ است که $\fn{Un}(x,x)$ و~$g(y)$ هر دو تعریف
شده باشند.

توجه کنید برای~$x$ ثابت، اگر $\cfind{x}(x)$ تعریف‌نشده باشد،
$h(x,y)$ برای هر~$y$ تعریف‌نشده است؛ و اگر $\cfind{x}(x)$ تعریف شده
باشد، $h(x,y)\simeq g(y)$. پس به ازای هر مقدار ثابت~$x$، $h(x,y)$
یا درست مانند $f$ رفتار می‌کند یا درست مانند $g$؛ و تصمیم دربارهٔ
تعریف‌شدن $\cfind{x}(x)$ همان تصمیم دربارهٔ این است که کدام‌یک از
این دو حالت برقرار است. اما این نیز همان تصمیم دربارهٔ عضویت یا عدم
عضویت $h_x(y)\simeq h(x,y)$ در~$C$ است؛ و اگر $A$ محاسبه‌پذیر بود،
دقیقاً می‌توانستیم چنین تصمیمی بگیریم.

به‌طور رسمی‌تر، چون $h$ محاسبه‌پذیر جزئی و دوموضعی است، برای نمایه‌ای چون~$e$
برابر تابع $\cfind{e}[2]$ است. بنا بر قضیهٔ $s$-$m$-$n$ تابع بازگشتی
اولیه‌ای چون~$s$ وجود دارد که به ازای هر~$x$،
$\cfind{s(e,x)}(y)\simeq h_x(y)$. اکنون برای هر $x$، اگر
$\cfind{x}(x)\fdefined$، آنگاه $\cfind{s(e,x)}$ همان تابع~$g$ است و
پس $s(e,x)\in A$. از سوی دیگر، اگر $\cfind{x}(x)\uparrow$، آنگاه
$\cfind{s(e,x)}$ همان تابع~$f$ است و پس $s(e,x)\notin A$. به بیان
دیگر، برای هر~$x$، $x\in K$ اگر و تنها اگر $s(e,x)\in A$. اگر $A$
محاسبه‌پذیر بود، $K$ نیز چنین می‌شد که تناقض است. پس $A$
محاسبه‌پذیر نیست.
\end{proof}

قضیهٔ رایس بسیار نیرومند است. نتیجهٔ فوری زیر چند کاربرد نمونه را
نشان می‌دهد.
\begin{cor}
مجموعه‌های زیر تصمیم‌ناپذیرند.
\begin{enumerate}
\item $\Setabs{x}{\text{$17$ در برد $\cfind{x}$ است}}$
\item $\Setabs{x}{\text{$\cfind{x}$ ثابت است}}$
\item $\Setabs{x}{\text{$\cfind{x}$ تام است}}$
\item $\Setabs{x}{\text{هرگاه $y < y'$، $\cfind{x}(y) \fdefined$، و
    اگر $\cfind{x}(y') \fdefined$، آنگاه $\cfind{x}(y) < \cfind{x}(y')$}}$
\end{enumerate}
\end{cor}

\begin{proof}
  همهٔ این‌ها مجموعه‌های نمایه‌های نابدیهی‌اند.
\end{proof}

\end{document}
